\section{Reprojection Error e Golden Rule}
\label{sec:golden_rule}

%─────────────────────────────────────────────
\subsection{Reprojection error}
\label{sub:reproj}

Una volta stimati \(\mat{K}\), \(\mathbf{k}_r=(k_1,k_2,p_1,p_2,k_3)\)
e gli estrinseci \(\{[\mat{R}_j|\mathbf{t}_j]\}\), la qualità della
calibrazione si misura con il \emph{reprojection error} medio:
\[
  \epsilon_{\mathrm{RMSE}}
  =\sqrt{\frac{1}{N}\sum_{j=1}^{n}\sum_{i=1}^{m_j}
          \bigl\|\homo{x}_{ji}-\hat{\homo{x}}_{ji}\bigr\|^2},
\]
dove \(\hat{\homo{x}}_{ji}=\pi\!\left(\mat{K}[\mat{R}_j|\mathbf{t}_j]
\homo{X}_i,\mathbf{k}_r\right)\) è la riproiezione del punto 3D.

%─────────────────────────────────────────────
\subsection{La regola aurea: proiezione su una varietà}
\label{sub:golden_rule}

I parametri stimati nella fase lineare definiscono un punto
\(\hat\theta\) nello spazio dei parametri, ma i dati misurati con
rumore producono un sistema \emph{inconsistente}: non esiste un
\(\theta\) che soddisfi esattamente tutte le equazioni.
La \emph{Gold Standard algorithm} afferma che la soluzione ottimale non
è quella che azzera le equazioni algebriche, ma quella che minimizza la
\textbf{distanza geometrica} tra i punti misurati \(\tilde{P}\) e la
superficie \(V\) delle configurazioni ammissibili:
\[
  \hat{\theta}^* = \arg\min_{\theta}\sum_i
  d\!\bigl(\tilde{P}_i,\,V(\theta)\bigr)^2.
\]
Geometricamente, il punto ottimo \(P_0\in V\) è il punto della
superficie più vicino a \(\tilde{P}\): la distanza minima si raggiunge
quando il vettore \((\tilde{P}-P_0)\) è ortogonale al piano tangente
a \(V\) in \(P_0\), ovvero quando
\[
  \nabla f(P_0)^T(\tilde{P}-P_0) = 0.
\]
Il gradiente \(\nabla f(P_0)\) individua la direzione normale alla
superficie nel punto di minima distanza.

\begin{figure}[H]
\centering
\begin{tikzpicture}[>=Stealth, scale=1.3]

  % ── Etichetta spazio dei parametri (spostata a sinistra) ─────
  \node[font=\small, gray!50!black] at (-2.55, 1.1)
    {spazio dei parametri \(\theta\)};

  % ── Superficie V (curva) ─────────────────────────────────────
  \draw[blue!70!black, very thick, domain=-1.6:2.4, smooth, samples=60]
    plot (\x, {0.18*\x*\x - 0.12*\x - 0.5})
    node[right, font=\small] {\(V(\theta)\)};

  % ── Punto P0 sulla superficie ────────────────────────────────
  \coordinate (P0) at (0.8, -0.4808);
  \fill[blue!80!black] (P0) circle (2.2pt);
  \node[below right, font=\small, blue!80!black] at (P0) {\(P_0\)};

  % ── Punto misurato P tilde ───────────────────────────────────
  \coordinate (Pt) at (1.55, 0.95);
  \fill[red!80!black] (Pt) circle (2.2pt);
  \node[above right, font=\small, red!80!black] at (Pt)
    {\(\tilde{P}\)};

  % ── Distanza geometrica (perpendicolare) ─────────────────────
  \draw[red!70!black, dashed, thick] (Pt) -- (P0)
    node[midway, right, font=\small, red!60!black]
    {\(d(\tilde{P},V)\)};

  % ── Piano tangente in P0 ─────────────────────────────────────
  \draw[green!50!black, thick]
    ($(P0) + (-1.1, -0.185)$) -- ($(P0) + (1.1, 0.185)$)
    node[right, font=\small, green!40!black] {\small tangente};

  % ── Gradiente in P0 ──────────────────────────────────────────
  \draw[->, orange!80!black, very thick]
    (P0) -- ($(P0) + (-0.25, 1.48)$)
    node[left, font=\small, orange!70!black]
    {\(\nabla f(P_0)\)};

  % ── Angolo retto ─────────────────────────────────────────────
  \draw[gray!60]
    ($(P0) + (0.10, 0.60)$) --
    ($(P0) + (0.00, 0.617)$) --
    ($(P0) + (-0.00, 0.017)$);

\end{tikzpicture}
\caption{Regola aurea: il punto ottimo \(P_0\in V(\theta)\) è il punto
         della superficie ammissibile più vicino a \(\tilde{P}\). Il
         vettore \((\tilde{P}-P_0)\) è ortogonale al piano tangente,
         ovvero parallelo al gradiente \(\nabla f(P_0)\).}
\label{fig:golden_rule}
\end{figure}

%─────────────────────────────────────────────
\subsection{Applicazione al progetto: DLT, SVD e LM}
\label{sub:pipeline_calibrazione}

Il problema della regola d'oro applicato alla calibrazione è che
la superficie delle configurazioni ammissibili \(V\) è
\emph{non lineare}: i parametri di \(\mat{K}\), \(\mat{R}_j\),
\(\mathbf{t}_j\) e \(\mathbf{k}_r\) entrano nella funzione di
riproiezione in modo polinomiale, quindi non è possibile trovare
\(\hat\theta^*\) in forma chiusa. La pipeline adottata risolve il
problema in due fasi:

\begin{enumerate}

  \item \textbf{Inizializzazione lineare (DLT + SVD).}
  Per ogni frame \(j\) si costruisce il sistema lineare omogeneo
  \(\mat{A}_j\vect{m}=\vect{0}\) tramite \emph{Direct Linear
  Transform} (DLT): le corrispondenze punto--immagine vengono scritte
  come equazioni lineari negli elementi di \(\mat{M}_j\).
  Il sistema è sovradeterminato e si risolve ai minimi quadrati
  tramite \emph{Singular Value Decomposition} (SVD): la soluzione è
  il vettore singolare destro corrispondente al valore singolare
  minimo di \(\mat{A}_j\), che individua la direzione nello spazio
  dei parametri più vicina al nucleo. Dalle omografie \(\mat{M}_j\)
  si ricava \(\mat{K}\) per fattorizzazione di Cholesky su
  \(\mat{\Gamma}^{-1}\). Il risultato \(\hat\theta_0\) è una buona
  approssimazione iniziale, ma non giace ancora su \(V\) poiché
  ignora distorsione e rumore.

  \item \textbf{Raffinamento non lineare (Levenberg--Marquardt).}
  Partendo da \(\hat\theta_0\), l'algoritmo di
  \emph{Levenberg--Marquardt} (LM) si muove iterativamente nello
  spazio dei parametri verso il minimo di \(\epsilon_{\mathrm{RMSE}}\).
  Ad ogni passo combina la robustezza della \emph{discesa del
  gradiente} (quando è lontano dal minimo) con la convergenza rapida
  del metodo di \emph{Gauss--Newton} (quando è vicino), proiettando
  progressivamente \(\hat\theta\) sulla superficie ammissibile \(V\).

\end{enumerate}
